\documentclass[12pt,a4paper]{article}
\usepackage[utf-8]{inputenc}
\usepackage[portuguese]{babel}
\usepackage[margin=2.5cm]{geometry}
\usepackage{amsmath}
\usepackage{amssymb}
\usepackage{listings}
\usepackage{xcolor}
\usepackage{hyperref}
\usepackage{array}
\usepackage{booktabs}
\usepackage{tikz}
\usepackage{graphicx}
\usepackage{fancyhdr}
\usepackage{lastpage}

% Configuração de cores
\definecolor{codebackground}{rgb}{0.95,0.95,0.95}
\definecolor{darkred}{rgb}{0.8,0,0}

% Estilo de código
\lstset{
    language=Java,
    backgroundcolor=\color{codebackground},
    basicstyle=\ttfamily\small,
    keywordstyle=\color{darkred}\bfseries,
    commentstyle=\color{gray},
    stringstyle=\color{darkred},
    breaklines=true,
    tabsize=2,
    showstringspaces=false,
    numbers=left,
    numberstyle=\tiny
}

% Headers e footers
\pagestyle{fancy}
\fancyhf{}
\rhead{GraphNet Analyzer}
\lhead{Documentação Técnica}
\cfoot{\thepage~de~\pageref{LastPage}}

% Título
\title{\textbf{GraphNet Analyzer}\\
\large Documentação Técnica Detalhada\\
\normalsize Análise Arquitetural e de Dependências}

\author{
    Jafte Carneiro Fagundes da Silva\\
    Nicolas Hrescak\\
    \\
    \textit{Orientador: Prof. Fabrício Enembreck}\\
    Pontifícia Universidade Católica do Paraná (PUCPR)\\
    Bacharelado em Ciência da Computação
}

\date{16 de junho de 2026}

\begin{document}

\maketitle
\thispagestyle{empty}
\newpage

\tableofcontents
\newpage

% ================================================================
% SEÇÃO 1: INTRODUÇÃO
% ================================================================
\section{Introdução}

GraphNet Analyzer é uma aplicação completa de análise e visualização de grafos de alta dimensionalidade, desenvolvida em Java puro (JDK 8+), combinando uma interface gráfica moderna com um sistema de linha de comando legado. O projeto implementa uma rede social com 5.000 usuários e mais de 25.000 conexões direcionadas.

\subsection{Objetivos do Projeto}

\begin{enumerate}
    \item Implementar uma estrutura de dados de grafo eficiente com suporte a vértices nomeados e arestas ponderadas
    \item Implementar 7 algoritmos clássicos de análise de grafos (conectividade, componentes, ciclos, Euleriano, centralidades)
    \item Fornecer uma interface gráfica interativa com visualização em tempo real
    \item Implementar um sistema de recomendação de seguidores via busca em profundidade (DFS)
    \item Aplicar princípios SOLID para obter código modular, testável e extensível
\end{enumerate}

\subsection{Escopo Técnico}

\begin{itemize}
    \item \textbf{Linguagem:} Java 8+ (sem dependências externas)
    \item \textbf{GUI:} Swing com componentes customizados
    \item \textbf{Tamanho:} 5.000 vértices, ~25.000 a 40.000 arestas
    \item \textbf{Formato de I/O:} Pajek (.net)
    \item \textbf{Padrões de Design:} MVC, Facade, Factory, Strategy, Dependency Injection
\end{itemize}

\newpage

% ================================================================
% SEÇÃO 2: ARQUITETURA GERAL
% ================================================================
\section{Arquitetura Geral do Sistema}

\subsection{Visão de Módulos}

A aplicação está organizada em 6 camadas funcionais, cada uma com responsabilidade bem definida:

\begin{table}[h]
\centering
\begin{tabular}{|l|l|l|}
\hline
\textbf{Camada} & \textbf{Pacote} & \textbf{Responsabilidade} \\
\hline
\textbf{Apresentação} & \texttt{graph.gui} & Interface gráfica e componentes customizados \\
\hline
\textbf{Controladores} & \texttt{graph.gui.*Handler} & Roteamento de eventos e orquestração \\
\hline
\textbf{Domínio} & \texttt{graph.domain} & Lógica de negócio (recomendações) \\
\hline
\textbf{Algoritmos} & \texttt{graph.algorithm} & Análise de grafos (7 algoritmos) \\
\hline
\textbf{Estrutura de Dados} & \texttt{graph.model} & Representação interna do grafo \\
\hline
\textbf{I/O e Utilitários} & \texttt{graph.io}, \texttt{graph.generator} & Importação, exportação e geração \\
\hline
\end{tabular}
\end{table}

\subsection{Diagrama de Dependências entre Pacotes}

As dependências entre pacotes seguem a regra de acoplamento mínimo:

\begin{verbatim}
    Main, Menu
         ↓
    GraphGUI ← (orquestração)
         ↓
    Handlers (Loading, Algorithm, SocialNetwork)
         ↓
    GraphUIState (estado centralizado)
         ↓
    Graph ← (lê)
         ↓
    GraphAlgorithms, SocialNetwork
         ↓
    GraphGenerator, PajekIO, NamesLoader
\end{verbatim}

A estrutura garante que:
\begin{itemize}
    \item \textbf{Algoritmos} nunca dependem de GUI
    \item \textbf{Domínio} nunca conhece Handlers
    \item \textbf{Graph} é imutável para leitura segura
    \item \textbf{Handlers} recebem dependências via injeção de construtor
\end{itemize}

\newpage

% ================================================================
% SEÇÃO 3: CAMADA DE ESTRUTURA DE DADOS
% ================================================================
\section{Camada de Estrutura de Dados}

\subsection{Classe \texttt{Graph}}

A classe \texttt{Graph} é a estrutura de dados central, representando um grafo ponderado, rotulado e opcionalmente direcionado.

\subsubsection{Responsabilidades}

\begin{itemize}
    \item Armazenar vértices em um array de strings (\texttt{names[]})
    \item Armazenar arestas em lista de adjacências (\texttt{adj[]})
    \item Fornecer API pública para consulta segura de vizinhos
    \item Gerenciar automaticamente inversas de arestas para grafos não-direcionados
\end{itemize}

\subsubsection{Atributos Principais}

\begin{table}[h]
\centering
\begin{tabular}{|l|l|l|}
\hline
\textbf{Atributo} & \textbf{Tipo} & \textbf{Descrição} \\
\hline
\texttt{numVertices} & \texttt{final int} & Número total de vértices \\
\hline
\texttt{directed} & \texttt{boolean} & Se o grafo é direcionado \\
\hline
\texttt{names[]} & \texttt{private String[]} & Rótulos dos vértices \\
\hline
\texttt{adj[]} & \texttt{private List<int[]>[]} & Lista de adjacências (destino, peso) \\
\hline
\end{tabular}
\end{table}

\subsubsection{Métodos Públicos}

\begin{table}[h]
\centering
\begin{tabular}{|l|l|}
\hline
\textbf{Método} & \textbf{Descrição e Complexidade} \\
\hline
\texttt{Graph(int, boolean)} & Construtor. O(V) \\
\hline
\texttt{setName(int, String)} & Define rótulo do vértice. O(1) \\
\hline
\texttt{getName(int)} & Retorna rótulo. O(1) \\
\hline
\texttt{addEdge(int, int, int)} & Adiciona aresta. O(1) amortizado \\
\hline
\texttt{findByName(String)} & Busca vértice por nome (case-insensitive). O(V) \\
\hline
\texttt{neighbors(int)} & Retorna lista de vizinhos. O(1) \\
\hline
\end{tabular}
\end{table}

\subsubsection{Fluxo de Criação}

\begin{lstlisting}[language=Java, caption=Uso da classe Graph]
// Criação
Graph graph = new Graph(5000, true); // direcionado

// Configuração de nomes
for (int i = 0; i < 5000; i++) {
    graph.setName(i, nomes[i]);
}

// Adição de arestas
graph.addEdge(0, 1, 1); // pessoa 0 segue pessoa 1
\end{lstlisting}

\newpage

% ================================================================
% SEÇÃO 4: CAMADA DE ALGORITMOS
% ================================================================
\section{Camada de Algoritmos}

\subsection{Classe \texttt{GraphAlgorithms}}

Implementa 7 algoritmos clássicos de análise de grafos, especializados para grafos direcionados com componentes fracamente conectados.

\subsubsection{Responsabilidades}

\begin{itemize}
    \item Verificar conectividade (fraca para direcionados)
    \item Identificar componentes conexos
    \item Detectar ciclos (coloração DFS)
    \item Verificar propriedades Eulerianas
    \item Calcular centralidade de proximidade (Dijkstra)
    \item Calcular centralidade de intermediação (Algoritmo de Brandes)
\end{itemize}

\subsubsection{Algoritmo 1: Conectividade}

\begin{verbatim}
Função: isConnected()
Entrada: Grafo G
Saída: boolean

1. Se |V| ≤ 1, retorna true
2. Cria vetor visitados[|V|] ← false
3. Inicia BFS a partir do vértice 0 no grafo não-direcionado
4. Retorna true se todos os vértices foram visitados
\end{verbatim}

\textbf{Complexidade:} O(V + E) \\
\textbf{Caso de Uso:} Verificar se a rede social é conexa (fracamente)

\subsubsection{Algoritmo 2: Componentes Conexos}

\begin{verbatim}
Função: getComponents()
Entrada: Grafo G
Saída: List<List<Integer>>

1. Cria vetor visitados[|V|] ← false
2. Para cada vértice i não visitado:
   a. Cria nova lista de componentes
   b. Inicia BFS a partir de i
   c. Marca todos os vértices alcançáveis como visitados
   d. Adiciona lista de componentes ao resultado
3. Retorna lista de componentes
\end{verbatim}

\textbf{Complexidade:} O(V + E) \\
\textbf{Caso de Uso:} Identificar subgrupos desconectados na rede

\subsubsection{Algoritmo 3: Euleriano}

\begin{verbatim}
Função: isEulerian()
Entrada: Grafo G
Saída: 2 (Ciclo) | 1 (Caminho) | 0 (Não Euleriano)

Para grafos NÃO-direcionados:
  1. Conta vértices com grau ímpar
  2. Se 0 → Ciclo Euleriano (retorna 2)
  3. Se 2 → Caminho Euleriano (retorna 1)
  4. Senão → Não Euleriano (retorna 0)

Para grafos DIRECIONADOS:
  1. Conta out-degree e in-degree de cada vértice
  2. Se todos iguais → Ciclo Euleriano (retorna 2)
  3. Se exatamente um com out-degree = in-degree + 1
     e outro com in-degree = out-degree + 1 → Caminho Euleriano (retorna 1)
  4. Senão → Não Euleriano (retorna 0)
\end{verbatim}

\textbf{Complexidade:} O(V + E) \\
\textbf{Caso de Uso:} Verificar se existe caminho/ciclo que visite toda aresta exatamente uma vez

\subsubsection{Algoritmo 4: Detecção de Ciclo}

\begin{verbatim}
Função: hasCycle()
Entrada: Grafo G
Saída: boolean

Para grafos DIRECIONADOS (coloração DFS):
  1. Cores: BRANCO (0, não visitado), CINZA (1, visitando), PRETO (2, visitado)
  2. Para cada vértice não visitado:
     a. Inicia DFS com cor CINZA
     b. Se encontrar vizinho CINZA → ciclo detectado, retorna true
     c. Marca como PRETO ao terminar
  3. Retorna false se nenhum ciclo encontrado

Para grafos NÃO-direcionados (rastreamento de pai):
  1. Para cada vértice não visitado:
     a. Inicia DFS
     b. Se encontrar vizinho visitado (não é pai) → ciclo, retorna true
  3. Retorna false se nenhum ciclo encontrado
\end{verbatim}

\textbf{Complexidade:} O(V + E) \\
\textbf{Caso de Uso:} Detectar influências circulares na rede social

\subsubsection{Algoritmo 5: Centralidade de Proximidade}

\begin{verbatim}
Função: closenessCentrality()
Entrada: Grafo G ponderado
Saída: double[] closeness

Para cada vértice s:
  1. Executa Dijkstra a partir de s
  2. Calcula soma de distâncias para todos os outros vértices
  3. closeness[s] = (|V| - 1) / soma_distâncias

Retorna array closeness
\end{verbatim}

\textbf{Complexidade:} O(V × (V + E) log V) com Dijkstra e heap binário \\
\textbf{Caso de Uso:} Identificar usuários centrais na rede (fácil acesso a outros)

\subsubsection{Algoritmo 6: Centralidade de Intermediação (Brandes)}

\begin{verbatim}
Função: betweennessCentrality()
Entrada: Grafo G ponderado
Saída: double[] betweenness

Para cada vértice s:
  1. Executa Dijkstra a partir de s
  2. Constrói lista de predecessores para cada vértice
  3. Retrocede do vértice com maior distância até s
  4. Acumula dependências: δ(s,t) = |predecessores(t)| × δ(s,predecessores(t)) / |predecessores(s)|
  5. Adiciona δ(s,t) à betweenness[t]

Divide betweenness por (V-1)×(V-2)/2 para grafos não-direcionados
Retorna array betweenness
\end{verbatim}

\textbf{Complexidade:} O(V × (V + E) log V) \\
\textbf{Caso de Uso:} Identificar usuários que atuam como "ponte" entre grupos

\subsubsection{Algoritmo 7: Dijkstra (Auxiliar)}

\begin{verbatim}
Função: dijkstra(int source)
Entrada: Vértice de origem
Saída: int[] distâncias

1. Inicializa distâncias com INFINITO, exceto source = 0
2. Cria PriorityQueue com (distância, vértice)
3. Enquanto queue não vazia:
   a. Extrai (d, u) com menor distância
   b. Se d > distância[u], continua (já processado)
   c. Para cada vizinho v de u:
      - Se distância[u] + peso(u,v) < distância[v]:
        • Atualiza distância[v]
        • Adiciona (nova_distância, v) à queue
4. Retorna array de distâncias
\end{verbatim}

\textbf{Complexidade:} O((V + E) log V) com heap binário \\
\textbf{Caso de Uso:} Encontrar caminhos mínimos em redes ponderadas

\subsubsection{Métodos Auxiliares Privados}

\begin{table}[h]
\centering
\begin{tabular}{|l|l|}
\hline
\textbf{Método} & \textbf{Descrição} \\
\hline
\texttt{bfsUndirected(boolean[])} & BFS ignorando direção das arestas \\
\hline
\texttt{undirectedNeighbors(int)} & Retorna vizinhos em ambas direções \\
\hline
\texttt{dfsColor(int, int[], int)} & DFS com coloração (branco/cinza/preto) \\
\hline
\texttt{dfsCycle(int, int, boolean[], boolean[])} & DFS para detectar ciclos (não-direcionado) \\
\hline
\texttt{dijkstra(int)} & Dijkstra para uma fonte \\
\hline
\texttt{backtrackBrandes(...)} & Retrocesso em Brandes \\
\hline
\end{tabular}
\end{table}

\newpage

% ================================================================
% SEÇÃO 5: CAMADA DE DOMÍNIO
% ================================================================
\section{Camada de Domínio}

\subsection{Classe \texttt{SocialNetwork}}

Implementa a lógica de negócio específica da rede social — o sistema de recomendação de seguidores.

\subsubsection{Responsabilidades}

\begin{itemize}
    \item Validar nomes de usuários no grafo
    \item Verificar relacionamentos diretos (já segue)
    \item Encontrar caminhos indiretos via DFS
    \item Gerar recomendações com caminho completo
\end{itemize}

\subsubsection{Fluxo do Método \texttt{recommendFollower}}

\begin{verbatim}
Função: recommendFollower(sourceName, targetName)
Entrada: Nomes de dois usuários
Saída: String (recomendação com caminho ou motivo da recusa)

1. Valida nomes: busca índices no grafo
   - Se qualquer um não encontrado → "Person(s) not found"

2. Rejeita auto-recomendação: u == v
   - Se igual → "Cannot recommend a person to themselves"

3. Verifica relacionamento direto: (u,v) já existe?
   - Se sim → "already follows"

4. Busca caminho indireto via DFS
   - Chama dfsPath(u, v, visited[], path[])
   - Se caminho encontrado:
     • Constrói string "Recommended! Path found:\n  u → ... → v"
   - Senão:
     • "Not recommended. No path exists..."

5. Retorna mensagem com resultado
\end{verbatim}

\subsubsection{Algoritmo DFS com Rastreamento de Caminho}

\begin{verbatim}
Função PRIVADA: dfsPath(current, dest, visited[], path[])
Entrada: Vértice atual, destino, visitados, caminho acumulado
Saída: boolean (verdadeiro se destino alcançado)

1. Marca current como visitado
2. Adiciona current ao final de path[]
3. Se current == dest, retorna true (encontrado!)
4. Para cada vizinho v de current não visitado:
   a. Recursivamente chama dfsPath(v, dest, ...)
   b. Se retornar true, retorna true (propaga sucesso)
5. Remove current do final de path[] (backtrack)
6. Retorna false (destino não alcançável daqui)
\end{verbatim}

\textbf{Complexidade:} O(V + E) no pior caso \\
\textbf{Propriedade:} Encontra o primeiro caminho (não necessariamente o mais curto)

\subsubsection{Exemplo de Uso}

\begin{lstlisting}[language=Java, caption=Sistema de Recomendação de Seguidores]
SocialNetwork network = new SocialNetwork(graph);

String result = network.recommendFollower("Alice", "Carol");
// Possíveis saídas:
// 1. "Recommended! Path found: Alice -> Bob -> Carol"
// 2. "'Alice' already follows 'Carol'."
// 3. "Not recommended. No path exists..."
// 4. "Person(s) not found in the graph."
\end{lstlisting}

\newpage

% ================================================================
% SEÇÃO 6: CAMADA DE I/O E UTILITÁRIOS
% ================================================================
\section{Camada de I/O e Utilitários}

\subsection{Classe \texttt{GraphGenerator}}

Responsável pela criação de grafos segundo diferentes estratégias.

\subsubsection{Métodos Estáticos}

\begin{table}[h]
\centering
\begin{tabular}{|l|l|}
\hline
\textbf{Método} & \textbf{Descrição} \\
\hline
\texttt{generateSocialNetwork(String namesFile)} & Gera rede social com 5.000 pessoas \\
\hline
\texttt{generateRandomGraph(int V, int E, boolean)} & Gera grafo aleatório \\
\hline
\end{tabular}
\end{table}

\subsubsection{Geração de Rede Social}

\begin{verbatim}
Função: generateSocialNetwork(String namesFile)
Entrada: Caminho para arquivo de nomes
Saída: Graph (direcionado, 5000 vértices)

1. Carrega 5.000 nomes do arquivo via NamesLoader
2. Cria grafo direcionado com 5000 vértices
3. Atribui nomes a cada vértice
4. Para cada pessoa i:
   a. Escolhe aleatoriamente entre 5 e 8 pessoas para seguir
   b. Arestas têm peso 1 (representa "seguir")
5. Resultado: ~25.000 a 40.000 arestas
\end{verbatim}

\subsection{Classe \texttt{NamesLoader}}

Utilitário para carregamento do arquivo de nomes.

\begin{table}[h]
\centering
\begin{tabular}{|l|l|}
\hline
\textbf{Método} & \textbf{Descrição} \\
\hline
\texttt{loadNames(String)} & Retorna String[] com nomes do arquivo \\
\hline
\end{tabular}
\end{table}

\subsection{Classe \texttt{PajekIO}}

Implementa importação e exportação no formato Pajek (.net).

\subsubsection{Formato Pajek}

\begin{verbatim}
*Vertices <count>
<id> "<name>"
...
*Arcs
<from> <to> <weight>
...
\end{verbatim}

\subsubsection{Métodos Estáticos}

\begin{table}[h]
\centering
\begin{tabular}{|l|l|}
\hline
\textbf{Método} & \textbf{Descrição} \\
\hline
\texttt{importFrom(String filename)} & Lê arquivo em pajek/input/ \\
\hline
\texttt{export(Graph, String filename)} & Escreve em pajek/output/ \\
\hline
\end{tabular}
\end{table}

\newpage

% ================================================================
% SEÇÃO 7: CAMADA DE APRESENTAÇÃO (GUI)
% ================================================================
\section{Camada de Apresentação}

\subsection{Classe \texttt{GraphGUI}}

Orquestrador principal da interface gráfica, responsável pela estrutura da janela e roteamento de eventos.

\subsubsection{Responsabilidades}

\begin{itemize}
    \item Criar estrutura de janela (BorderLayout)
    \item Inicializar componentes (Header, Sidebar esquerdo, GraphPanel, Sidebar direito)
    \item Instanciar handlers e injetar dependências
    \item Roteador de eventos para handlers especializados
    \item Atualizar UI com resultados
\end{itemize}

\subsubsection{Arquitetura Refatorada}

A classe foi refatorada de 1600+ linhas para ~250 linhas de orquestração pura:

\begin{table}[h]
\centering
\begin{tabular}{|l|l|}
\hline
\textbf{Responsabilidade Original} & \textbf{Extraída Para} \\
\hline
Carregamento/geração de grafos & GraphLoadingHandler \\
\hline
Execução de algoritmos & AlgorithmHandler \\
\hline
Operações de rede social & SocialNetworkHandler \\
\hline
Gerenciamento de estado & GraphUIState \\
\hline
Renderização de grafo & GraphPanel \\
\hline
Design tokens & Theme \\
\hline
\end{tabular}
\end{table}

\subsubsection{Componentes da GUI}

\begin{verbatim}
┌─────────────────────────────────────────────────┐
│             GraphGUI (JFrame)                   │
├─────────────────────────────────────────────────┤
│  Header: Título, status (V, E, conexo?)         │
├─────────────────────────────────────────────────┤
│ │                                             │ │
│ │  Sidebar Esquerdo          GraphPanel      │ Sidebar Direito
│ │  - Gerar                   (Renderização   │ - Estatísticas
│ │  - Importar                do Grafo)       │ - Top Nós
│ │  - Algoritmos                              │ - Resultados
│ │                                            │
│ │                                            │
└─────────────────────────────────────────────────┘
\end{verbatim}

\subsection{Classe \texttt{GraphLoadingHandler}}

Controller para operações de carregamento e geração de grafos.

\subsubsection{Responsabilidades}

\begin{itemize}
    \item Coordenar geração de rede social
    \item Coordenar geração de grafo aleatório
    \item Coordenar importação de arquivo Pajek
    \item Atualizar estado da aplicação
    \item Chamar callbacks de progresso
\end{itemize}

\subsubsection{Fluxo de Carregamento}

\begin{verbatim}
Usuário clica "Gerar Rede Social"
       ↓
GraphLoadingHandler.onGenerateSocialNetwork()
       ↓
GraphGenerator.generateSocialNetwork("names.txt")
       ↓
GraphUIState.setGraph(novoGrafo)
       ↓
GraphGUI.updateUI() / GraphPanel.render()
       ↓
Resultados aparecem na tela
\end{verbatim}

\subsection{Classe \texttt{AlgorithmHandler}}

Controller para execução de algoritmos de análise.

\subsubsection{Responsabilidades}

\begin{itemize}
    \item Validar tamanho do grafo
    \item Executar algoritmos de GraphAlgorithms
    \item Formatar resultados para exibição
    \item Alertar usuário sobre operações pesadas
    \item Atualizar lista de top nós (centralidades)
\end{itemize}

\subsubsection{Métodos Principais}

\begin{table}[h]
\centering
\begin{tabular}{|l|l|}
\hline
\textbf{Método} & \textbf{Descrição} \\
\hline
\texttt{checkConnectivity()} & Chama algs.isConnected() \\
\hline
\texttt{getComponents()} & Chama algs.getComponents() \\
\hline
\texttt{checkEulerian()} & Chama algs.isEulerian() \\
\hline
\texttt{hasCycle()} & Chama algs.hasCycle() \\
\hline
\texttt{closenessCentrality()} & Chama algs.closenessCentrality() \\
\hline
\texttt{betweennessCentrality()} & Chama algs.betweennessCentrality() \\
\hline
\end{tabular}
\end{table}

\subsection{Classe \texttt{SocialNetworkHandler}}

Controller para operações de rede social.

\subsubsection{Responsabilidades}

\begin{itemize}
    \item Processar requisições de recomendação
    \item Validar entrada do usuário
    \item Chamar SocialNetwork.recommendFollower()
    \item Formatar caminho encontrado para exibição
    \item Atualizar resultado na interface
\end{itemize}

\subsection{Classe \texttt{GraphPanel}}

Renderizador visual do grafo em um canvas (JPanel).

\subsubsection{Responsabilidades}

\begin{itemize}
    \item Desenhar vértices (círculos coloridos)
    \item Desenhar arestas (setas direcionadas)
    \item Suportar pan (movimento) interativo
    \item Suportar zoom (ampliação/redução)
    \item Usar design tokens de Theme
\end{itemize}

\subsubsection{Algoritmo de Renderização}

\begin{verbatim}
paintComponent(Graphics2D g):
  1. Limpa canvas
  2. Para cada aresta (u, v):
     - Calcula posição de u e v usando layout (circular ou força)
     - Desenha linha com seta direcionada
  3. Para cada vértice u:
     - Desenha círculo colorido
     - Desenha rótulo centralizado
  4. Aplica transformações (pan + zoom)
\end{verbatim}

\subsection{Classe \texttt{GraphUIState}}

Gerenciador centralizado de estado da aplicação.

\subsubsection{Estado Gerenciado}

\begin{table}[h]
\centering
\begin{tabular}{|l|l|}
\hline
\textbf{Propriedade} & \textbf{Tipo} \\
\hline
\texttt{currentGraph} & \texttt{Graph} \\
\hline
\texttt{algorithms} & \texttt{GraphAlgorithms} \\
\hline
\texttt{socialNetwork} & \texttt{SocialNetwork} \\
\hline
\texttt{lastResults} & \texttt{String} \\
\hline
\texttt{topNodes} & \texttt{List<Integer>} \\
\hline
\end{tabular}
\end{table}

\subsection{Classe \texttt{Theme}}

Sistema de design tokens centralizado.

\subsubsection{Tokens Fornecidos}

\begin{table}[h]
\centering
\begin{tabular}{|l|l|}
\hline
\textbf{Token} & \textbf{Valor} \\
\hline
\texttt{PRIMARY\_COLOR} & Cor primária da aplicação \\
\hline
\texttt{SECONDARY\_COLOR} & Cor secundária \\
\hline
\texttt{BACKGROUND\_COLOR} & Cor de fundo \\
\hline
\texttt{FONT\_FAMILY} & Font padrão \\
\hline
\texttt{FONT\_SIZE} & Tamanho de fonte padrão \\
\hline
\texttt{PADDING} & Espaçamento padrão \\
\hline
\end{tabular}
\end{table}

\subsection{Componentes Customizados}

\begin{itemize}
    \item \textbf{PrimaryButton} — Botão primário com hover effects
    \item \textbf{MenuActionButton} — Botão com ícone e label
    \item \textbf{CollapsiblePanel} — Painel expansível/colapsável
    \item \textbf{ProgressDialog} — Diálogo com barra de progresso
    \item \textbf{VectorIcon} — Ícones vetoriais escaláveis
\end{itemize}

\newpage

% ================================================================
% SEÇÃO 8: FLUXOS DE DADOS PRINCIPAIS
% ================================================================
\section{Fluxos de Dados Principais}

\subsection{Fluxo 1: Inicialização da Aplicação}

\begin{verbatim}
Main.main(String[])
  ↓
Menu.chooseConsoleMode()
  ├─ Se GUI (padrão):
  │   └─ SwingUtilities.invokeLater()
  │        └─ new GraphGUI().setVisible(true)
  │              ├─ initComponents()
  │              ├─ createHandlers()
  │              ├─ setupListeners()
  │              └─ setVisible(true)
  │
  └─ Se Console:
      └─ Main.showMainMenu()
           └─ Loop interativo CLI
\end{verbatim}

\subsection{Fluxo 2: Geração de Rede Social}

\begin{verbatim}
Usuário clica "Gerar Rede Social"
  ↓
GraphGUI.onGenerateSocialNetwork() [ActionListener]
  ↓
GraphLoadingHandler.generateSocialNetwork()
  ├─ Mostra ProgressDialog
  ├─ GraphGenerator.generateSocialNetwork("names.txt")
  │   ├─ NamesLoader.loadNames("names.txt")
  │   ├─ Graph g = new Graph(5000, true)
  │   ├─ setName(i, nomes[i]) para cada pessoa
  │   └─ addEdge() para cada relação "segue"
  ├─ GraphUIState.setGraph(g)
  │   └─ Recria GraphAlgorithms e SocialNetwork
  ├─ GraphPanel.setGraph(g)
  └─ updateResults()
       └─ Exibe estatísticas no sidebar direito
\end{verbatim}

\subsection{Fluxo 3: Execução de Algoritmo (Exemplo: Conectividade)}

\begin{verbatim}
Usuário clica "Conectividade"
  ↓
GraphGUI.onCheckConnectivity() [ActionListener]
  ↓
AlgorithmHandler.checkConnectivity()
  ├─ GraphAlgorithms algs = state.getAlgorithms()
  ├─ boolean connected = algs.isConnected()
  │   ├─ Se |V| ≤ 1, retorna true
  │   ├─ BFS no grafo não-direcionado
  │   └─ Verifica se todos vértices foram visitados
  ├─ Formata resultado em String
  └─ GraphGUI.updateResults(resultado)
       └─ Exibe "Connected: Yes/No" no sidebar direito
\end{verbatim}

\subsection{Fluxo 4: Sistema de Recomendação de Seguidores}

\begin{verbatim}
Usuário digita "Alice" e "Carol"
  ↓
GraphGUI.onRecommendFollower()
  ↓
SocialNetworkHandler.recommendFollower("Alice", "Carol")
  ├─ SocialNetwork social = state.getSocialNetwork()
  ├─ String resultado = social.recommendFollower("Alice", "Carol")
  │   ├─ Busca índices: u = findByName("Alice"), v = findByName("Carol")
  │   ├─ Valida (não nulos, não iguais)
  │   ├─ Verifica se u já segue v diretamente
  │   ├─ Se não, chama dfsPath(u, v, visitados[], caminho[])
  │   │   ├─ DFS percorre grafo
  │   │   ├─ Se encontrar v, reconstrói caminho
  │   │   └─ Retorna List<Integer> com sequência de vértices
  │   └─ Formata resultado: "Recommended! Path: Alice → ... → Carol"
  ├─ GraphGUI.updateResults(resultado)
  └─ GraphPanel destaca caminho (opcional)
\end{verbatim}

\subsection{Fluxo 5: Centralidade de Intermediação (Brandes)}

\begin{verbatim}
Usuário clica "Betweenness Centrality"
  ↓
AlgorithmHandler.betweennessCentrality()
  ├─ Alerta: "Large graph (5000 vertices). Continue? (y/n)"
  ├─ GraphAlgorithms algs = state.getAlgorithms()
  ├─ double[] b = algs.betweennessCentrality()
  │   └─ Para cada vértice s (0 a 4999):
  │        ├─ dijkstra(s) → calcula distâncias de s
  │        ├─ Reconstrói árvore de caminhos mínimos
  │        └─ Acumula dependências em betweenness[]
  ├─ Ordena vértices por betweenness descrescente
  ├─ Extrai top 10
  └─ GraphGUI.updateResults() e lstTopNodes
       └─ Exibe nomes e valores no sidebar direito
\end{verbatim}

\newpage

% ================================================================
% SEÇÃO 9: PADRÕES DE DESIGN
% ================================================================
\section{Padrões de Design}

\subsection{MVC (Model-View-Controller)}

\begin{table}[h]
\centering
\begin{tabular}{|l|l|l|}
\hline
\textbf{Componente} & \textbf{Papel} & \textbf{Classe(s)} \\
\hline
\textbf{Model} & Dados e lógica & Graph, GraphAlgorithms, SocialNetwork \\
\hline
\textbf{View} & Apresentação & GraphGUI, GraphPanel, componentes customizados \\
\hline
\textbf{Controller} & Orquestração & Handlers (Loading, Algorithm, SocialNetwork) \\
\hline
\end{tabular}
\end{table}

\subsection{Facade}

\texttt{GraphGUI} expõe interface unificada para manipulação de handlers:

\begin{lstlisting}[language=Java, caption=Facade em GraphGUI]
public void onGenerateSocialNetwork() {
    loadingHandler.generateSocialNetwork();
}

public void onCheckConnectivity() {
    algorithmHandler.checkConnectivity();
}

public void onRecommendFollower(String src, String tgt) {
    socialHandler.recommendFollower(src, tgt);
}
\end{lstlisting}

\subsection{Dependency Injection}

Handlers recebem dependências via construtor:

\begin{lstlisting}[language=Java, caption=Injeção de Dependência]
// Em GraphGUI.initComponents():
GraphUIState state = new GraphUIState();
GraphLoadingHandler handler =
    new GraphLoadingHandler(state, graphPanel, progressCallback);
\end{lstlisting}

\subsection{Factory}

\texttt{GraphGenerator} usa factory methods estáticos:

\begin{lstlisting}[language=Java, caption=Factory Pattern]
Graph g1 = GraphGenerator.generateSocialNetwork("names.txt");
Graph g2 = GraphGenerator.generateRandomGraph(100, 500, true);
\end{lstlisting}

\subsection{Strategy}

Múltiplos algoritmos em \texttt{GraphAlgorithms}:

\begin{lstlisting}[language=Java, caption=Strategy Pattern]
double[] closeness = algs.closenessCentrality();   // Dijkstra
double[] between = algs.betweennessCentrality();   // Brandes
boolean hasCycle = algs.hasCycle();                // DFS Coloração
\end{lstlisting}

\subsection{Observer}

Callbacks para atualização de UI:

\begin{lstlisting}[language=Java, caption=Observer Pattern]
loadingHandler.setProgressCallback((percent) -> {
    progressDialog.setValue(percent);
    SwingUtilities.invokeLater(() -> repaint());
});
\end{lstlisting}

\newpage

% ================================================================
% SEÇÃO 10: ANÁLISE DE DEPENDÊNCIAS
% ================================================================
\section{Análise de Dependências}

\subsection{Matriz de Dependências entre Classes}

\begin{table}[h]
\centering
\begin{tabular}{|l|l|l|l|l|}
\hline
\textbf{Classe} & \textbf{Depende De} & \textbf{Usado Por} & \textbf{Tipo} \\
\hline
\texttt{Graph} & Nenhuma & Todos & Estrutura \\
\hline
\texttt{GraphAlgorithms} & Graph & Handlers, Main & Algoritmos \\
\hline
\texttt{SocialNetwork} & Graph & Handlers, Main & Domínio \\
\hline
\texttt{GraphGenerator} & Graph, NamesLoader & Main, Handlers & Factory \\
\hline
\texttt{NamesLoader} & File I/O & GraphGenerator & Utilitário \\
\hline
\texttt{PajekIO} & Graph, File I/O & Main, Handlers & I/O \\
\hline
\texttt{GraphGUI} & Todos os Handlers & Main & Orquestrador \\
\hline
\texttt{GraphLoadingHandler} & GraphUIState, GraphPanel & GraphGUI & Controller \\
\hline
\texttt{AlgorithmHandler} & GraphUIState & GraphGUI & Controller \\
\hline
\texttt{SocialNetworkHandler} & GraphUIState & GraphGUI & Controller \\
\hline
\texttt{GraphUIState} & Graph, Algs, Social & Handlers & Estado \\
\hline
\texttt{GraphPanel} & Graph & GraphGUI, Handlers & Vista \\
\hline
\texttt{Theme} & Nenhuma & Componentes & Tokens \\
\hline
\end{tabular}
\end{table}

\subsection{Direção das Dependências}

\begin{verbatim}
GUI (não depende de backend)
  ↓
Handlers (orquestram operações)
  ↓
GraphUIState (gerencia estado)
  ↓
Graph ← GraphAlgorithms (não-direcionado, usa o mesmo Graph)
        ← SocialNetwork (não-direcionado, usa o mesmo Graph)
  ↓
GraphGenerator, PajekIO, NamesLoader (utilitários para criar/carregar)
\end{verbatim}

\subsection{Garantias de Acoplamento Baixo}

\begin{itemize}
    \item \textbf{Algoritmos são puros:} GraphAlgorithms não conhece GUI
    \item \textbf{Domínio é puro:} SocialNetwork não conhece Handlers
    \item \textbf{Graph é read-only:} Não há callbacks ou listeners
    \item \textbf{Handlers recebem dependências:} Não instanciam suas dependências
    \item \textbf{Inversão de controle:} Main/GraphGUI orquestram, não as classes
\end{itemize}

\newpage

% ================================================================
% SEÇÃO 11: ANÁLISE DE COMPLEXIDADE
% ================================================================
\section{Análise de Complexidade}

\subsection{Algoritmos}

\begin{table}[h]
\centering
\begin{tabular}{|l|l|l|l|}
\hline
\textbf{Algoritmo} & \textbf{Tempo} & \textbf{Espaço} & \textbf{Observação} \\
\hline
Conectividade & O(V+E) & O(V) & BFS \\
\hline
Componentes & O(V+E) & O(V) & BFS \\
\hline
Euleriano & O(V+E) & O(V) & Contagem de graus \\
\hline
Ciclo & O(V+E) & O(V) & DFS com coloração \\
\hline
Dijkstra & O((V+E)logV) & O(V) & Heap binário \\
\hline
Closeness & O(V(V+E)logV) & O(V) & V×Dijkstra \\
\hline
Betweenness (Brandes) & O(V(V+E)logV) & O(V) & V×Dijkstra \\
\hline
\end{tabular}
\end{table}

\subsection{Complexidade para Rede Social (V=5000, E=30000)}

\begin{table}[h]
\centering
\begin{tabular}{|l|l|l|}
\hline
\textbf{Operação} & \textbf{Complexidade Teórica} & \textbf{Tempo Estimado} \\
\hline
Conectividade & O(35000) & <1 segundo \\
\hline
Componentes & O(35000) & <1 segundo \\
\hline
Euleriano & O(35000) & <1 segundo \\
\hline
Ciclo & O(35000) & <1 segundo \\
\hline
Closeness (Top 10) & O(5000×35000×log5000) & 5-10 minutos \\
\hline
Betweenness (Top 10) & O(5000×35000×log5000) & 5-10 minutos \\
\hline
Recomendação (DFS) & O(35000) no pior caso & <1 segundo típico \\
\hline
\end{tabular}
\end{table}

Nota: Centralidades são operações pesadas. A aplicação alerta usuário antes de executar em grafos > 500 vértices.

\newpage

% ================================================================
% SEÇÃO 12: ESTRUTURA DE PACOTES E RESPONSABILIDADES
% ================================================================
\section{Estrutura de Pacotes e Responsabilidades}

\subsection{Pacote \texttt{graph.model}}

\begin{table}[h]
\centering
\begin{tabular}{|l|l|}
\hline
\textbf{Classe} & \textbf{Responsabilidade} \\
\hline
\texttt{Graph} & Estrutura de dados de grafo \\
\hline
\end{tabular}
\end{table}

\subsection{Pacote \texttt{graph.algorithm}}

\begin{table}[h]
\centering
\begin{tabular}{|l|l|}
\hline
\textbf{Classe} & \textbf{Responsabilidade} \\
\hline
\texttt{GraphAlgorithms} & 7 algoritmos + helpers privados \\
\hline
\end{tabular}
\end{table}

\subsection{Pacote \texttt{graph.domain}}

\begin{table}[h]
\centering
\begin{tabular}{|l|l|}
\hline
\textbf{Classe} & \textbf{Responsabilidade} \\
\hline
\texttt{SocialNetwork} & Lógica de recomendação + DFS \\
\hline
\end{tabular}
\end{table}

\subsection{Pacote \texttt{graph.io}}

\begin{table}[h]
\centering
\begin{tabular}{|l|l|}
\hline
\textbf{Classe} & \textbf{Responsabilidade} \\
\hline
\texttt{PajekIO} & Importação/exportação (estático) \\
\hline
\texttt{NamesLoader} & Carregamento de nomes (estático) \\
\hline
\end{tabular}
\end{table}

\subsection{Pacote \texttt{graph.generator}}

\begin{table}[h]
\centering
\begin{tabular}{|l|l|}
\hline
\textbf{Classe} & \textbf{Responsabilidade} \\
\hline
\texttt{GraphGenerator} & Factory para criação de grafos (estático) \\
\hline
\end{tabular}
\end{table}

\subsection{Pacote \texttt{graph.gui}}

\begin{table}[h]
\centering
\begin{tabular}{|l|l|}
\hline
\textbf{Classe} & \textbf{Responsabilidade} \\
\hline
\texttt{GraphGUI} & Orquestração de UI + roteamento de eventos \\
\hline
\texttt{GraphPanel} & Renderização de grafo \\
\hline
\texttt{GraphUIState} & Gerenciamento centralizado de estado \\
\hline
\texttt{GraphLoadingHandler} & Carregamento/geração \\
\hline
\texttt{AlgorithmHandler} & Execução de algoritmos \\
\hline
\texttt{SocialNetworkHandler} & Recomendações \\
\hline
\texttt{ProgressDialog} & Diálogo de progresso \\
\hline
\texttt{Theme} & Design tokens \\
\hline
\texttt{PrimaryButton} & Botão customizado \\
\hline
\texttt{MenuActionButton} & Botão de ação customizado \\
\hline
\texttt{CollapsiblePanel} & Painel expansível \\
\hline
\texttt{VectorIcon} & Ícones vetoriais \\
\hline
\end{tabular}
\end{table}

\subsection{Pacote \texttt{graph} (raiz)}

\begin{table}[h]
\centering
\begin{tabular}{|l|l|}
\hline
\textbf{Classe} & \textbf{Responsabilidade} \\
\hline
\texttt{Main} & Ponto de entrada + menu CLI \\
\hline
\texttt{Menu} & Seletor GUI/CLI \\
\hline
\end{tabular}
\end{table}

\newpage

% ================================================================
% SEÇÃO 13: PRINCÍPIOS SOLID APLICADOS
% ================================================================
\section{Princípios SOLID Aplicados}

\subsection{S — Single Responsibility Principle}

Cada classe tem uma única responsabilidade bem definida:

\begin{itemize}
    \item \texttt{Graph} — Apenas armazenar e consultar estrutura
    \item \texttt{GraphAlgorithms} — Apenas executar algoritmos
    \item \texttt{SocialNetwork} — Apenas lógica de domínio
    \item \texttt{GraphGUI} — Apenas orquestração de UI
    \item \texttt{Handlers} — Apenas roteamento de eventos
\end{itemize}

\subsection{O — Open/Closed Principle}

Estrutura aberta para extensão, fechada para modificação:

\begin{itemize}
    \item Novos algoritmos podem ser adicionados a \texttt{GraphAlgorithms} sem modificar código existente
    \item Novos componentes customizados podem ser criados herdando de \texttt{JComponent}
    \item Novo handlers podem ser criados herdando interface padrão
\end{itemize}

\subsection{L — Liskov Substitution Principle}

Componentes customizados substituem \texttt{JComponent} sem quebrar contrato:

\begin{lstlisting}[language=Java, caption=LSP em Componentes Customizados]
JComponent btn = new PrimaryButton("Click me");
JComponent panel = new CollapsiblePanel("Title");
// Ambos podem ser usados onde JComponent é esperado
\end{lstlisting}

\subsection{I — Interface Segregation Principle}

Interfaces pequenas e focadas:

\begin{lstlisting}[language=Java, caption=ISP em Listeners]
loadingHandler.setProgressCallback(value -> {...});
// Não forçamos handler a implementar interface monolítica
\end{lstlisting}

\subsection{D — Dependency Inversion Principle}

Dependências apontam para abstrações, não para concretos:

\begin{lstlisting}[language=Java, caption=DIP em Handlers]
// Handler não cria GraphUIState, recebe via construtor
handler = new AlgorithmHandler(state);
// Desacopla criação de injeção
\end{lstlisting}

\newpage

% ================================================================
% SEÇÃO 14: REFACTORIZAÇÕES REALIZADAS
% ================================================================
\section{Refactorizações Realizadas}

\subsection{Refatoração 1: Extração de Algoritmos}

\textbf{Antes:} Graph era uma God Class com ~340 linhas contendo algoritmos, geração e I/O.

\textbf{Depois:}
\begin{itemize}
    \item Algoritmos extraídos para \texttt{GraphAlgorithms} (classe dedicada)
    \item Geração extraída para \texttt{GraphGenerator}
    \item I/O extraído para \texttt{PajekIO}
    \item Graph reduzida para ~100 linhas de estrutura pura
\end{itemize}

\subsection{Refatoração 2: Refactoring de GUI}

\textbf{Antes:} GraphGUI era monolítica com 1600+ linhas.

\textbf{Depois:}
\begin{itemize}
    \item Orquestração pura em GraphGUI (~250 linhas)
    \item Carregamento em GraphLoadingHandler
    \item Algoritmos em AlgorithmHandler
    \item Domínio em SocialNetworkHandler
    \item Estado em GraphUIState
    \item Renderização em GraphPanel
    \item Design tokens em Theme
\end{itemize}

\subsection{Refatoração 3: Encapsulamento}

\textbf{Antes:} Campo \texttt{adj} era package-private, permitindo acesso direto.

\textbf{Depois:}
\begin{itemize}
    \item \texttt{adj} agora é \texttt{private}
    \item Novo método público \texttt{neighbors(int)} para acesso seguro
    \item GraphAlgorithms e SocialNetwork usam \texttt{neighbors()} em vez de acesso direto
\end{itemize}

\subsection{Refatoração 4: Eliminação de Duplicação}

\textbf{Antes:} Código de geração de nomes era duplicado.

\textbf{Depois:}
\begin{itemize}
    \item Classe dedicada \texttt{NamesLoader}
    \item \texttt{GraphGenerator} chama \texttt{NamesLoader}
    \item Uma única fonte de verdade para carregamento de nomes
\end{itemize}

\newpage

% ================================================================
% SEÇÃO 15: CONCLUSÃO
% ================================================================
\section{Conclusão}

\subsection{Arquitetura Alcançada}

GraphNet Analyzer demonstra uma arquitetura bem estruturada, seguindo os princípios SOLID e padrões de design clássicos. A separação clara entre camadas (Apresentação, Domínio, Algoritmos, Estrutura de Dados) resulta em código:

\begin{itemize}
    \item \textbf{Modular:} Cada classe tem responsabilidade bem definida
    \item \textbf{Testável:} Componentes podem ser testados isoladamente
    \item \textbf{Extensível:} Novos algoritmos/componentes podem ser adicionados
    \item \textbf{Manutenível:} Baixo acoplamento facilita manutenção
    \item \textbf{Reusável:} Classes podem ser usadas em outros projetos
\end{itemize}

\subsection{Características Principais}

\begin{enumerate}
    \item \textbf{Estrutura de Dados Eficiente:} Lista de adjacências para O(1) de adição de arestas
    \item \textbf{7 Algoritmos Clássicos:} Cobertura completa de análise de grafos
    \item \textbf{Interface Gráfica Moderna:} Visualização interativa com componentes customizados
    \item \textbf{Sistema de Recomendação:} DFS com rastreamento de caminho completo
    \item \textbf{Múltiplos Modos:} GUI e CLI para diferentes casos de uso
    \item \textbf{I/O Flexível:} Formato Pajek para interoperabilidade
\end{enumerate}

\subsection{Autoria}

Este projeto foi desenvolvido colaborativamente por:

\begin{itemize}
    \item \textbf{Jafte Carneiro Fagundes da Silva} (@cyberfika)
    \item \textbf{Nicolas Hrescak} (@NicolasHrescak)
\end{itemize}

Sob orientação do Prof. Fabrício Enembreck, PUCPR — Pontifícia Universidade Católica do Paraná.

Todos os arquivos de código foram criados pelos autores, sem uso de templates ou geradores de código automático. A documentação interna (Javadoc) está em 100\% de cobertura em português.

\subsection{Recomendações para Expansão Futura}

\begin{enumerate}
    \item \textbf{Layout de Visualização:} Implementar Force-Directed Layout (Fruchterman-Reingold) para melhor visualização de grafos grandes
    \item \textbf{Persistência:} Adicionar suporte a banco de dados para grafos muito grandes (>50k vértices)
    \item \textbf{Testes Automatizados:} Implementar suite de testes JUnit para aumentar confiabilidade
    \item \textbf{Paralelização:} Usar ExecutorService para paralelizar Dijkstra/Brandes
    \item \textbf{APIs Adicionais:} Expor como REST API para consumo por outras aplicações
\end{enumerate}

\appendix

\section{Referências de Código}

\subsection{Arquivo de Nomes}

Localização: \texttt{data/names.txt}

Formato: Um nome brasileiro por linha (5.000 nomes)

\subsection{Formato Pajek}

Arquivo de exemplo: \texttt{pajek/input/example.net}

\begin{verbatim}
*Vertices 3
1 "Alice"
2 "Bob"
3 "Carol"
*Arcs
1 2 1
2 3 1
1 3 1
\end{verbatim}

\subsection{Comando de Compilação}

\begin{verbatim}
javac -encoding UTF-8 -d out \
  src/graph/model/Graph.java \
  src/graph/algorithm/GraphAlgorithms.java \
  src/graph/domain/SocialNetwork.java \
  src/graph/io/PajekIO.java \
  src/graph/io/NamesLoader.java \
  src/graph/generator/GraphGenerator.java \
  src/graph/gui/Theme.java \
  src/graph/gui/VectorIcon.java \
  src/graph/gui/GraphUIState.java \
  src/graph/gui/ProgressDialog.java \
  src/graph/gui/PrimaryButton.java \
  src/graph/gui/MenuActionButton.java \
  src/graph/gui/CollapsiblePanel.java \
  src/graph/gui/GraphPanel.java \
  src/graph/gui/GraphLoadingHandler.java \
  src/graph/gui/AlgorithmHandler.java \
  src/graph/gui/SocialNetworkHandler.java \
  src/graph/gui/GraphGUI.java \
  src/graph/Menu.java \
  src/graph/Main.java
\end{verbatim}

\subsection{Comando de Execução}

\begin{verbatim}
java -cp out graph.Main
\end{verbatim}

Ao iniciar, escolha:
\begin{itemize}
    \item [1] Interface Gráfica (padrão, recomendado)
    \item [2] Modo Console (legado)
\end{itemize}

\end{document}
